% El contenido de dudas.md, de la raíz del repositorio, para mostrarlo impreso. Si
% cambia uno, cambia el otro. El PDF se copia a la raíz como dudas.pdf.
\documentclass[12pt,a4paper]{article}
% Preámbulo del informe de la Etapa 3: el de la Etapa 2 sin la notación de la
% gramática. Computer Modern a 12pt, como los documentos de la cátedra, y todo en
% negro.
\usepackage[T1]{fontenc}
\usepackage[utf8]{inputenc}
\usepackage{lmodern}
\usepackage{amsmath,amssymb}
\usepackage[spanish,es-noquoting,es-noshorthands,es-tabla]{babel}
\usepackage[a4paper,margin=2.6cm,headsep=14pt]{geometry}
\usepackage{microtype}
\usepackage{booktabs}
\usepackage{longtable}
\usepackage{array}
\usepackage{xcolor}
\usepackage{fancyhdr}
\usepackage{titlesec}
\usepackage{enumitem}
\usepackage{fancyvrb}
\usepackage[hidelinks,bookmarks=true]{hyperref}

\usepackage{float}
\usepackage{graphicx}
\usepackage{pdflscape}
\DeclareUnicodeCharacter{2014}{---}
\DeclareUnicodeCharacter{2026}{\ldots}

% --- Paleta ------------------------------------------------------------------
\definecolor{tinta}{rgb}{0,0,0}

\color{tinta}

% --- Títulos ------------------------------------------------------------------
\titleformat{\section}{\Large\bfseries\color{tinta}}{\thesection}{0.7em}{}
\titleformat{\subsection}{\large\bfseries\color{tinta}}{\thesubsection}{0.7em}{}
\titleformat{\subsubsection}{\bfseries\color{tinta}}{\thesubsubsection}{0.7em}{}
\titlespacing*{\section}{0pt}{2.4ex plus 1ex minus .2ex}{1.3ex plus .2ex}

% --- Encabezado y pie ---------------------------------------------------------
\setlength{\headheight}{14pt}
\pagestyle{fancy}
\fancyhf{}
\renewcommand{\headrulewidth}{0.4pt}
\renewcommand{\headrule}{\hbox to\headwidth{\color{tinta}\leaders\hrule height \headrulewidth\hfill}}
\fancyhead[L]{\footnotesize\color{tinta}Compiladores e Intérpretes — Etapa 3}
\fancyhead[R]{\footnotesize\color{tinta}\nouppercase{\leftmark}}
\fancyfoot[C]{\footnotesize\color{tinta}\thepage}

% --- Bloques de consola --------------------------------------------------------
% \footnotesize y no \small: la línea más larga de una sesión de consola tiene
% 82 caracteres y a \small no entra en la caja de texto.
\DefineVerbatimEnvironment{consola}{Verbatim}%
  {fontsize=\footnotesize,xleftmargin=1.1em,formatcom=\color{tinta}}

\setlength{\emergencystretch}{4.5em}
\tolerance=3000

\setlist[itemize]{leftmargin=1.3em,itemsep=0.25ex,topsep=0.5ex}
\setlist[enumerate]{leftmargin=1.5em,itemsep=0.25ex,topsep=0.5ex}
\setlist[description]{leftmargin=1.3em,itemsep=0.35ex,topsep=0.6ex}
\renewcommand{\arraystretch}{1.25}

\DeclareUnicodeCharacter{2192}{\ensuremath{\rightarrow}}
\fancyhead[R]{\footnotesize\color{tinta}Dudas}

\newcommand{\cod}[1]{\texttt{#1}}
\newcommand{\pregunta}[1]{\par\medskip\needspace{4\baselineskip}\noindent\textbf{#1}\par\smallskip}
\DefineVerbatimEnvironment{fuente}{Verbatim}%
  {fontsize=\small,xleftmargin=1.2em,frame=leftline,framerule=1.2pt,framesep=0.8em}
\usepackage{needspace}
\titleformat{\section}{\Large\bfseries\raggedright}{\thesection}{0.7em}{}
\setlength{\parindent}{0pt}
\setlength{\parskip}{0.55ex}

\begin{document}
\thispagestyle{fancy}
\raggedright

{\LARGE\bfseries Dudas --- Etapa 3 (chequeo de declaraciones)\par}
\smallskip
{\color{tinta}\rule{\textwidth}{0.6pt}}
\medskip

Cada punto trae un ejemplo, la salida que da hoy el compilador y por qué. Para contestar alcanza
con confirmar o con elegir la alternativa. Las líneas de los ejemplos cuentan desde 1.

\section*{Criterio general}

Usted indicó que los logros son independientes entre etapas. Por eso el semántico se implementa
para el MiniJava base del documento «Semántica de MiniJava – Parte I». Lo que agregan los logros
de la Etapa 2 (\cod{private}, genéricos anidados, atributos inicializados) se sigue aceptando
sintácticamente, pero no suma controles en esta etapa.

\textbf{¿Es correcto?}

\section*{A. Etapa 3}

\pregunta{1. ¿Hay que exigir un \cod{main}?}

\cod{class Aa \{ \}} → \cod{[SinErrores]}.

El documento no lo pide, y Java compila sin \cod{main}. La duda aparece porque los 5 casos
correctos del tester traen \cod{class Init \{ static void main() \{ \} \}}. Si hay que exigirlo,
falta saber dos cosas: si alcanza con que haya uno, y qué lexema y qué línea se reportan, porque
no hay ningún token al que apuntar.

\pregunta{2. \cod{private} y herencia.}

\begin{fuente}
class Aa { private int x; }
class Bb extends Aa { int x; }
\end{fuente}

→ \cod{[Error:x|2]}: el atributo ya está declarado en un ancestro.

Por el criterio general, la visibilidad no afecta ningún control. En Java este programa compila,
porque un miembro \cod{private} no se hereda.

\pregunta{3. Genéricos anidados.}

\begin{fuente}
class Caja<T> { }
class Par<T> { }
class Aa { Caja<Par<Nada>> x; }
\end{fuente}

→ \cod{[SinErrores]}, aunque \cod{Nada} no existe.

Se controla sólo el primer nivel: que \cod{Caja} sea genérica y que \cod{Par} exista. El documento
define el argumento como «el nombre de una clase o interfaz existente». Controlar el anidamiento es
el logro Genericidad Avanzada E3, que no perseguimos.

\pregunta{4. Herencia circular: ¿qué token se reporta?}

\begin{fuente}
class Aa extends Bb { }
class Bb extends Aa { }
\end{fuente}

→ \cod{[Error:Aa|1]}: la primera clase declarada que forma parte del ciclo.

Para lo que no tiene un token claramente asociado, la Pautas slide 30 pide usar el nombre de la
entidad con el problema. \cod{javac} marca el mismo lugar: «cyclic inheritance involving Aa».\\
Alternativa: \cod{[Error:Bb|1]}, el tipo que aparece en el \cod{extends}.

\pregunta{5. Argumento genérico sobre un tipo que no es genérico: ¿qué token se reporta?}

\begin{fuente}
class Caja { }
class Aa { Caja<String> x; }
\end{fuente}

→ \cod{[Error:Caja|2]}: el tipo que no declara parámetro. \cod{javac} dice «type Caja does not take
parameters».\\
Alternativa: \cod{[Error:String|2]}, el argumento.

\pregunta{6. \cod{extends String} y \cod{extends System}.}

\cod{class Aa extends String \{ \}} → \cod{[Error:String|1]}.

Es como en Java, donde las dos clases son \cod{final}: \cod{javac} dice «cannot inherit from final
String». El documento sólo dice que las dos «heredan de Object».

\pregunta{7. El parámetro de tipo en un método estático.}

\cod{class Caja<T> \{ static T crear() \{ \} \}} → \cod{[Error:T|1]}.

Es como en Java: \cod{javac} dice «non-static type variable T cannot be referenced from a static
context». El documento permite usar \cod{T} «en atributos, retornos y parámetros formales», y no
distingue los métodos estáticos.

\pregunta{8. Una interfaz no puede contradecir a \cod{Object}.}

\cod{interface Ix \{ int toString(); \}} → \cod{[Error:toString|1]}.

Es como en Java, donde una interfaz incluye los métodos públicos de \cod{Object}: \cod{javac} dice
«toString() in Ix cannot override toString() in Object».

\pregunta{9. Un método estático que coincide con uno de la interfaz: ¿qué token se reporta?}

\begin{fuente}
interface Ix { void m(); }
class Aa implements Ix { static void m() { } }
\end{fuente}

Que es error está fuera de duda: el documento §6.2 dice que un método estático no implementa uno
de interfaz.

→ \cod{[Error:m|2]}: el método que está mal. \cod{javac} marca el mismo lugar: «m() in Aa cannot
implement m() in Ix».\\
Alternativa: \cod{[Error:Aa|2]}, la clase que no implementa.

\pregunta{10. En esta etapa los cuerpos no se chequean.}

\cod{class Aa \{ void m() \{ this.b().a = 5; \} \}} → \cod{[SinErrores]}, aunque \cod{b} no está
declarado.

El documento §1 deja los cuerpos para la Etapa 4. Es el mismo caso que \cod{sintCorrecto04}, del
tester de la Etapa 2.

\section*{B. Para la Etapa 4 (lo que la Etapa 2 dejó al semántico)}

\pregunta{11. \cod{int x, y, z = 10;} (logro Variables Locales Clásicas).}

En Java esta línea inicializa \textbf{sólo} \cod{z}: si después se usa \cod{x}, \cod{javac} dice
«variable x might not have been initialized». ¿Se lee igual, o las tres variables valen 10?

\pregunta{12. Lo que el sintáctico deja pasar a propósito.}

Todos estos casos pasan hoy el análisis sintáctico. ¿Se controlan todos en la Etapa 4, como en
Java?

\smallskip
\begin{tabular}{>{\raggedright\arraybackslash}p{5.6cm} >{\raggedright\arraybackslash}p{9.4cm}}
\toprule
\textbf{Ejemplo} & \textbf{Qué se controla} \\
\midrule
\cod{5 = x;} & que el lado izquierdo del \cod{=} sea una variable \\
\cod{5 + 3;} & que la sentencia sea una asignación o una llamada (Java: «not a statement») \\
\cod{m()++;} & que el operando de \cod{++} o \cod{-{}-} sea una variable \\
\cod{noDeclarada = 5;} & que los nombres estén declarados \\
\cod{int x = "hola";} y \cod{return;} en un método \cod{int} & la compatibilidad de tipos \\
\cod{m(1, 2, 3)} con \cod{m} de un solo parámetro & la aridad de la llamada \\
\cod{int x; int x;} en un mismo bloque & las variables locales repetidas \\
\cod{new int[]\{\{\}\}} & que el inicializador coincida con las dimensiones \\
\cod{3.5e38f} & el rango del \cod{float} (Java: «floating-point number too large») \\
\cod{a + b * c} & la precedencia: la gramática de la cátedra pone los 13 operadores binarios en
un solo nivel. ¿Se aplica la de Java? \\
\bottomrule
\end{tabular}

\needspace{14\baselineskip}
\section*{C. Diferencias con los casos de prueba de un compañero}

Corrimos los 93 casos de prueba de un compañero y coincidimos en 86. En los 7 restantes, los dos
compiladores detectan el mismo error en la misma línea: sólo cambia el lexema reportado. Son
exactamente las consultas 4 y 5.

\pregunta{Herencia circular (3 casos).}

\begin{fuente}
class ClaseUno extends ClaseDos { }
class ClaseDos extends ClaseTres { }
class ClaseTres extends ClaseUno { }
\end{fuente}

→ Él reporta \cod{[Error:ClaseDos|1]}, el tipo del \cod{extends} de la primera clase. Nosotros,
\cod{[Error:ClaseUno|1]}, la clase.

Por qué: la Pautas slide 30 pide el nombre de la entidad con el problema cuando no hay un token
claramente asociado, y \cod{javac} marca la clase: «cyclic inheritance involving ClaseUno», con el
cursor en \cod{class ClaseUno}. Lo mismo pasa con un ciclo de interfaces y con uno de cuatro clases.

\pregunta{Argumento genérico sobre un tipo que no es genérico (4 casos).}

\begin{fuente}
class Simple { }
class Derivada extends Simple<String> { }
\end{fuente}

→ Él reporta \cod{[Error:String|2]}, el argumento. Nosotros, \cod{[Error:Simple|2]}, el tipo.

Por qué: \cod{String} existe y sería un argumento válido; lo que está mal es que \cod{Simple} no
declara un parámetro de tipo. \cod{javac} responsabiliza al tipo: «type Simple does not take
parameters». Lo mismo pasa con \cod{implements}, con un atributo y con un parámetro.

\end{document}
